\subsection{Interpolating and Folding ADCs}\label{sec:adc_interp_fold}

\subsubsection{Interpolating}
\paragraph{Overview:} flash preamplifiers (gain \(A_p\) around their threshold) at only every \(M\)-th
threshold; resistive dividers between neighbouring preamp outputs make the rest.\par
\noindent\textbf{Schematic:}\par
\begin{circuitfig}
  \draw (0,2) node[op amp, scale=0.7](Pa){}; \draw (0,-1) node[op amp, scale=0.7](Pb){};
  \node at (0,2.8) {$A_p$}; \node at (0,-1.8) {$A_p$};
  \draw (Pa.+) -- ++(-0.6,0) node[left]{$V_{IN}$}; \draw (Pa.-) -- ++(-0.6,0) node[left]{$V_{T,a}$};
  \draw (Pb.+) -- ++(-0.6,0) node[left]{$V_{IN}$}; \draw (Pb.-) -- ++(-0.6,0) node[left]{$V_{T,b}$};
  \draw (Pa.out) -- ++(0.6,0) coordinate(a) node[circ]{} node[above]{$V_a$};
  \draw (Pb.out) -- ++(0.6,0) coordinate(b) node[circ]{} node[below]{$V_b$};
  \draw (a) to[european resistor] ++(0,-1.5) node[circ]{} coordinate(j) to[european resistor] (b);
  \draw (j) -- ++(1.2,0) node[right]{$V_j \to$ comparator};
  \draw (a) -- ++(1.2,0) node[right]{$\to$ comparator};
  \draw (b) -- ++(1.2,0) node[right]{$\to$ comparator};
\end{circuitfig}
\paragraph{Transfer Function:}
\begin{enumerate}
  \item \textbf{Large Signal, Time Domain, DC Input Analysis (Q5)}: preamp outputs
    \(V_a = A_p(V_{IN} - V_{T,a})\), \(V_b = A_p(V_{IN} - V_{T,b})\); \(M\) equal resistors, tap \(j\) from \(V_a\)
    (superposition of the two ends):
\begin{align*}
  V_j &= \frac{(M - j) V_a + j V_b}{M} = \frac{(M - j) A_p(V_{IN} - V_{T,a}) + j A_p(V_{IN} - V_{T,b})}{M} \\
  V_j &= A_p\frac{(M - j + j) V_{IN} - (M - j) V_{T,a} - j V_{T,b}}{M} \\
  V_j &= A_p\left(V_{IN} - \frac{(M - j) V_{T,a} + j V_{T,b}}{M}\right), \hspace{1cm} j = 1 \ldots M - 1
\end{align*}
    zero crossing \(V_j = 0 \Rightarrow V_{IN} = V_{T,a} + \frac{j}{M}(V_{T,b} - V_{T,a})\): the linearly interpolated threshold.
  \item \textbf{Small Signal, Frequency Domain, AC Input Analysis (Q4)}: each preamp a single pole
    \(\omega_{pa}\), \(A_p(s) = \frac{A_p}{1 + s/\omega_{pa}}\); the divider is resistive, so every tap has the same pole.
\end{enumerate}
\paragraph{Analysis:}
\begin{enumerate}
  \item \textbf{Poles}: \(s_p = -\omega_{pa}\) at every tap (no extra delay mismatch between taps).
  \item \textbf{DC Gain}: \(\frac{dV_j}{dV_{IN}} = A_p\) at every tap (the \(\frac{M - j + j}{M} = 1\) of Q5).
  \item \textbf{Count}: preamps drop from \(2^n - 1\) to \(\approx \frac{2^n}{M}\); comparators remain \(2^n - 1\).
  \item \textbf{Impedances}: input capacitance falls by \(M\) (\(\frac{2^n}{M}\) preamps on \(V_{IN}\)).
  \item \textbf{Offset Averaging}: preamp offsets \(V_{os,a}\), \(V_{os,b}\) (independent, equal SD) shift the interpolated zero
    crossing by \(\frac{(M - j)V_{os,a} + jV_{os,b}}{M}\) [the Q5 tap formula]:
    \begin{align*}
      SD &= \frac{\sqrt{(M - j)^2 + j^2}}{M}SD(V_{os}) \leq SD(V_{os}) \hspace{1cm} [(M - j)^2 + j^2 \leq M^2 \text{ for } 0 \leq j \leq M]
    \end{align*}
    smallest at the middle tap (\(j = \frac{M}{2}\): \(\frac{SD(V_{os})}{\sqrt{2}}\)).
  \item \textbf{Advantages \& Limitations}: \(M\) times fewer preamps (input capacitance, power),
    averaging of preamp offsets; but valid only while the preamps are linear between neighbouring
    thresholds (\(|V_{IN} - V_T| < M X_{LSB}\)), and comparators remain \(2^n - 1\).
\end{enumerate}

\paragraph{Code:}
\begin{lstlisting}[language=Octave]
  clear; clc;
  pkg load symbolic;

  syms Vin VTa VTb Ap j
  syms M positive
  Vj = ((M - j)*Ap*(Vin - VTa) + j*Ap*(Vin - VTb))/M;
  printf('=== zero crossing of tap j ===\n'); pretty(simplify(solve(Vj == 0, Vin)))
  printf('=== offset SD factor at j = M/2 ===\n'); pretty(simplify(sqrt((M - M/2)^2 + (M/2)^2)/M))
\end{lstlisting}

\subsubsection{Folding}
\paragraph{Overview:} the input range mapped onto \(F\) folds; a coarse ADC finds the fold and a fine
ADC digitises the folded signal.\par
\noindent\textbf{Schematic:}\par
\begin{circuitfig}
  \draw[->] (-1,1) node[left]{$V_{IN}$} -- (0,1);
  \draw (0,1) node[circ]{} -- (0,2.2) -- (0.8,2.2); \draw (0,1) -- (0,-0.2) -- (0.8,-0.2);
  \draw (0.8,1.6) rectangle (3.2,2.8) node[midway, align=center]{coarse ADC\\ $\log_2 F$ bits};
  \draw (0.8,-0.8) rectangle (3.2,0.4) node[midway, align=center]{folding\\($F$ folds)};
  \draw[->] (3.2,-0.2) -- node[above]{$V_{fold}$} (4.6,-0.2);
  \draw (4.6,-0.8) rectangle (7,0.4) node[midway, align=center]{fine ADC\\ $\frac{2^n}{F} - 1$ comp.};
  \draw[->] (3.2,2.2) -- (7.8,2.2) |- (8,1);
  \draw[->] (7,-0.2) -- (7.8,-0.2) |- (8,1);
  \draw (8,0.4) rectangle (9.4,1.6) node[midway]{encode};
  \draw[->] (9.4,1) -- (10,1) node[right]{$D$};
\end{circuitfig}
\paragraph{Transfer Function:}
\begin{enumerate}
  \item \textbf{Large Signal, Time Domain, DC Input Analysis (Q5)}: in fold \(k\)
    \(V_{fold} = \pm\left(V_{IN} - k\frac{V_{REF}}{F}\right)\); coarse \(D_c = k\) (\(\log_2 F\) bits), fine
    \(D_f = \left\lfloor \frac{V_{fold}}{V_{REF}/F}\cdot\frac{2^n}{F} \right\rfloor\), \(D = D_c \frac{2^n}{F} + D_f\) (fine code reversed in odd folds).
  \item \textbf{Large Signal, Time Domain, AC Input Analysis (Q6)}: a full scale sine at \(f_{in}\) crosses
    all \(F\) folds every half period, so \(V_{fold}\) contains \(\approx F f_{in}\).
\end{enumerate}
\paragraph{Analysis:}
\begin{enumerate}
  \item \textbf{Comparator Count}: \(F - 1\) coarse plus \(\frac{2^n}{F} - 1\) fine,
\begin{align*}
  N(F) &= (F - 1) + \left(\frac{2^n}{F} - 1\right) \\
  \frac{dN}{dF} &= 1 - \frac{2^n}{F^2} = 0 \hspace{1cm} \Rightarrow \hspace{1cm} F = 2^{n/2} \\
  N_{\min} &= 2^{n/2} - 1 + \frac{2^n}{2^{n/2}} - 1 = 2 \cdot 2^{n/2} - 2
\end{align*}
    e.g. \(n = 8\): \(F = 2^4 = 16\), \(N = 2 \cdot 16 - 2 = 30\) instead of \(2^8 - 1 = 255\).
  \item \textbf{Bandwidth}: the folder must pass \(F f_{in}\) (Q6): \(f_{p,fold} > F f_{in,\max}\).
  \item \textbf{Advantages \& Limitations}: \(\approx 2 \cdot 2^{n/2}\) comparators at flash like speed; but the
    folder runs at \(F\) times the input frequency (bandwidth), fold corners are rounded (non-linear
    near the folds), and coarse/fine timing must align.
\end{enumerate}

\paragraph{Code:}
\begin{lstlisting}[language=Octave]
  clear; clc;
  pkg load symbolic;

  syms F n positive
  N = (F - 1) + (2^n/F - 1);
  Fopt = solve(diff(N, F) == 0, F);
  printf('=== optimum folds and comparators ===\n'); pretty(simplify(Fopt)); pretty(simplify(subs(N, F, Fopt)))
  printf('n = 8: F = %d, comparators = %d (flash %d)\n', 2^4, 2*16 - 2, 2^8 - 1);
\end{lstlisting}
